
% This LaTeX was auto-generated from an M-file by MATLAB.
% To make changes, update the M-file and republish this document.

\documentclass{article}
\usepackage{graphicx}
\usepackage{color}

\sloppy
\definecolor{lightgray}{gray}{0.5}
\setlength{\parindent}{0pt}

\begin{document}

    
    
\subsection*{Contents}

\begin{itemize}
\setlength{\itemsep}{-1ex}
   \item ES868 � Pr� Relat�rio Lab 3
   \item Controlador proporcional
   \item Controlador de Ziegler Nichols
\end{itemize}


\subsection*{ES868 � Pr� Relat�rio Lab 3}

\begin{par}
\textbf{Controle de plantas eletr�onicas utilizando um controlador PID digital}
\end{par} \vspace{1em}
\begin{par}
Turma A - 16/03/2015
\end{par} \vspace{1em}
\begin{itemize}
\setlength{\itemsep}{-1ex}
   \item Augusto Miranda Garcia - 104627
   \item Guilherme de Oliveira Souza � 117093
   \item Vin�cius Ragazi David - 120258
\end{itemize}
\begin{verbatim}
load ../G.mat
\end{verbatim}


\subsection*{Controlador proporcional}



\subsection*{Controlador de Ziegler Nichols}

\begin{par}
O primeiro passo no m�todo de Ziegler Nichols � encontrar $k = k_{osc}$ que deixa a planta no limiar de estabilidade quando em malha fechada:
\end{par} \vspace{1em}
\begin{par}
Para isso pode ser usado o crit�rio de Routh com a equa��o caracter�stica:
\end{par} \vspace{1em}
\begin{par}
$$1 + k G(s) = 0$$
\end{par} \vspace{1em}
\begin{par}
$$1 + \frac{8.993 k}{0.0003019 s^3 + 0.03622 s^2 + s} = 0$$
\end{par} \vspace{1em}
\begin{par}
$$\frac{0.0003019 s^3 + 0.03622 s^2 + s + 8.993 k}{0.0003019 s^3 + 0.03622 s^2 + s} = 0$$
\end{par} \vspace{1em}
\begin{par}
$$0.0003019 s^3 + 0.03622 s^2 + s + 8.993 k = 0$$
\end{par} \vspace{1em}
\begin{par}
Com o crit�rio de Routh
\end{par} \vspace{1em}



\end{document}
    
